\documentclass[11pt]{article}

% Packages
\usepackage[T1]{fontenc}
\usepackage[margin=1in]{geometry}
\usepackage{amsmath,amssymb}
\usepackage{natbib}
\usepackage{graphicx}
\usepackage{hyperref}
\graphicspath{{../../../figures/}}

\title{Draft Conclusions section}
\author{}
\date{}

\begin{document}
\maketitle

%% ================================================================
%% Rough draft of the Conclusions section for
%% 01_slr_forecast_intervention2026.tex. Standalone file, so any
%% \ref to a figure in the manuscript renders as ?? until pasted in.
%% ================================================================

\section{Conclusions}

SPECTRUM forecasts 21st-century sea-level rise from simple physical models of the individual contributors to sea level. We calibrate every contributor with a long observational record by Bayesian inference against that record. We take the others from published sensitivities. West Antarctic discharge is represented instead as a weighted mixture of scenarios, because the observational record cannot yet distinguish a self-sustaining retreat from one driven by ocean melting. Observed global mean sea level was withheld from every component calibration and reserved as a test of the assembled framework. Forced with observed temperatures, the summed components track global mean sea level from 1950 to 2025, a window spanning 1.1 to 1.2 $^\circ$C of warming. The hindcast rate over 1993 to 2025 is 3.19 mm\,yr$^{-1}$ against 3.17 to 3.35 mm\,yr$^{-1}$ in the observational records, and the hindcast reproduces the observed sensitivity of the rate to warming. It understates the rate over 1950 to 2000, where observational constraints on glaciers and the Greenland Ice Sheet are sparse.

Driven forward with warming trajectories, the framework projects medians and 90\% credible intervals at 2100 of 0.62 [0.44, 1.54] m at 2 $^\circ$C, 0.78 [0.60, 1.72] m at 3 $^\circ$C, and 0.99 [0.77, 1.91] m at 4 $^\circ$C, relative to the year 2000. Referenced to the AR6 baseline, these medians run 38 to 44\% above the IPCC AR6 medium-confidence projections \citep{Fox-Kemper2021}. Measured instead from pre-industrial sea level, those medians span 0.81 to 1.18 m, and the probability of exceeding 1 m is 26\% at 2 $^\circ$C, 44\% at 3 $^\circ$C, and 90\% at 4 $^\circ$C.

The uncertainty in these forecasts has two sources. The first is the set of contributors whose temperature sensitivity we calibrated against observations. These account for roughly 1\% of the forecast variance and 10 to 20\% of the width of the 90\% credible interval. Their contribution shifts systematically with warming, and emissions cuts reduce it. The second is the West Antarctic Ice Sheet. It accounts for nearly all of the forecast variance and for 80 to 90\% of the width of the 90\% credible interval. It responds weakly to the emissions pathway. Holding warming to 2 $^\circ$C rather than 4 $^\circ$C lowers the median by 0.37 m but moves the probability of exceeding 2 m only from 6\% to 3\%. Emissions policy sets how much rise is likely. West Antarctica sets how bad the worst case gets.

The two sources differ in what they cost and in what can be done about them. Every centimeter of rise from either source exposes roughly 1.6 million more people to annual coastal flooding, and about 90\% of them live in low- and middle-income countries \citep{Kulp2019, Jongman2012}. Resolving whether West Antarctica loses mass slowly or rapidly is worth on the order of US\$5T in avoided adaptation costs. That value is largest when confidence in a slow-WAIS future reaches 75\% and when the answer arrives early enough to shape infrastructure decisions. The answer will come from the rate of West Antarctic ice loss over the next two decades. Emissions cuts limit how far sea level rises. Sustained observation of West Antarctica tells coastal communities how much to prepare for. They need both.

%% ================================================================
\bibliographystyle{agufull08}
\bibliography{references}

\end{document}
